%----------------------------------------------------------------------------
%============================================================================
\section{Recovering fluids with known EoS in the EFT language}
%============================================================================
%-----------------------------------------------------------------------------

In this section we get some intuition for the form of the function $F(y,b)$ describing familiar fluids, such as cold dark matter and radiation.
Hopefully this will also reveal how the fluid EFT is more general than the usual scalar field description of fluids.
Recall the general expressions \eqref{eq:thermo_rho_p_nonbaro}-\eqref{eq:adiabatic_cs_nonadiabatic_pressure} for the general case and \eqref{eq:thermo_variables_barotropic} for the case $F=F(b)$, which is a particular instance of the barotropic case as discussed in Section \ref{sec:fluid_eft_exact}.

%========================================================================
\subsection{Fluids with constant EoS parameter}
%========================================================================

%-----------------------------------------------------------------------
\subsubsection{Dark Matter}
%-----------------------------------------------------------------------
%
Demanding Dark Matter obey $F=F(b)$ and imposing $w=0$ we get
\begin{align}
    0= w=\frac{bF_b-F}{F}\quad \Rightarrow\quad F=\Lambda^4 \,b\quad \text{for some scale } \, [\Lambda]=E.
\end{align}
This gives energy density $\rho=-F\propto \, b$ with the expecte behavior for cold matter $\rho \approx \, m \, n$.


%-------------------------------------------------------------------------
\subsubsection{Relativistic fluids}
%-------------------------------------------------------------------------
How to describe relativistic fluid in the EFT, like the photon-baryon fluid early on?

A fully relativistic fluid $p\simeq E$ that can be described in the Boltzmann formalism is necessarily barotropic, with the familiar EoS $w=\tfrac13$, namely
\begin{align}
    \rho:=T^0_0=\int\frac{d^3\vec{p}}{(2\pi)^3} E\, f, \quad P:= \tfrac13 T^i_i =\int\frac{d^3\vec{p}}{(2\pi)^3} \frac{p^2}{3E}= \frac13 \int\frac{d^3\vec{p}}{(2\pi)^3} E\, f = \tfrac13 \rho.
\end{align}
These formulas hold at any order since perturbations are encoded into $f=\bar{f}+\delta f$.
In particular we find that $c_a^2=\tfrac13$ necessarily, since
\begin{align}
    \delta P:=\int\frac{d^3\vec{p}}{(2\pi)^3} \frac{p^2}{3E}\, \delta f = \frac13 \int\frac{d^3\vec{p}}{(2\pi)^3} E \delta f =: \tfrac13 \delta\rho.
\end{align}
Similarly, there is no need to specify the thermodynamic variables we vary or keep fixed, since there is only one independent parameter be it $\rho,P,T,s,n,$.
For the record the entropy is
\begin{align}
    s=\frac{\rho+p}{T}= \frac{\tfrac43 \rho}{T}, \qquad \sigma=\frac{s}{n}=\left(\int\frac{d^3\vec{p}}{(2\pi)^3} f\right)^{-1} \frac{\tfrac43 \rho}{T}.
\end{align}

Consider now the EFt for the fluid.
Enforcing $F=F(b)$ and the relativistic EoS we get
\begin{align}
    \tfrac13 \overset{!}{=}w_{\textbf{baro}}=\frac{bF_b-F}{F}\quad \Rightarrow\quad F(b)=\Lambda^4\, b^{4/3}\,.
\end{align}
This captures the usual scaling for radiation $\rho\propto T^4 \propto n^{4/3}$.
Consistently, this automatically implies 
\begin{align}
    c_a^2=\frac{b^2F_{bb}}{bF_b}=\frac{1}{3}.
\end{align}

Are there fully relativistic fluids i.e. $w\equiv\tfrac13$ with $c_a^2\neq\tfrac13 ?$
No.
Indeed enforcing $w=\tfrac13$ in the fluid EFT implies the fluid is actually barotropic since $p=w\rho$ is a function of energy density only, and we only have one independent DoF in the thermodynamic plane.
The variables $y$ and $b$ are related via the constant equation of state 
\begin{align}
    \tfrac13=w=\frac{F-bF_b}{yF_y-F} \quad \Rightarrow\quad \tfrac43 F=\tfrac13 yF_y+bF_b.
\end{align}
In particular, differentiating this relation wrt $y$ and $b$ respectively and rearranging a bit we get
\begin{align}
    yF_y-ybF_{yb}=\tfrac13 y^2F_{yy}, \quad \tfrac13 (bF_b-ybF_{yb})=b^2F_{bb}.
\end{align}
Plugging these into the formulas \eqref{eq:adiabatic_cs_nonadiabatic_pressure} we get the expected behavior for a barotropic fluid
\begin{align}
    c_a^2=\frac{-y^2F_{yy}b^2F_{bb}+(yF_y-byF_{by})^2}{y^2F_{yy}(yF_y-bF_b)}\overset{identically}{=}\frac13\,, \quad \Gamma =0.
\end{align}

%-------------------------------------------------------------------------
\subsubsection{Fluids with arbitrary constant EoS parameter}
%-------------------------------------------------------------------------

This case is completely analogous to the relativistic fluid case, simply with general $w$ instead of $w=\tfrac13$.
The fluid is barotropic $p=w\rho$, and the variable $y$ can actually be out-integrated and expressed in terms of $b$ only with no further functional freedom.
Indeed from the equation of state we get a first-order PDE which can be solved by characteristics,
\begin{align}\label{eq:constant_eos_w}
    \frac{F-F_b}{yF_y-F}=w\approx \mathrm{const}\quad\Rightarrow\quad (1+w)F=w\, yF_y+bF_b.
\end{align}
Choosing convenitn variables $\gamma:=\log y$ and $\beta:=\log b$, the characteristics are
{\small
\begin{align}\begin{aligned}
    \dot{\beta}&=1\,,\\
    \dot{\gamma}&=w \quad \Rightarrow\quad y = y_0\, (b/b_0)^w, \\
    \dot{z}&=(1+w)z\,.
\end{aligned}\end{align}}
The general solution is, for an arbitrary function $f=f(z)$ on one variable,
\begin{align}\label{eq:constant_eos_w_solution_pde}
    F(y,b)=f(y\,b^{-w})\, b^{(1+w)}.
\end{align}
Note it reduces to the solution for $F=F(b)$ upon inserting $y = y_0\, (b/b_0)^w$, confirming that 
an EFT with $F=F(y,b)$ but $w=\text{const}.$ is can be reduced to $F=F(b)$.

The implications of $w=\mathrm{const}$ on $c_a^2$ and $\Gamma$ defined in \eqref{eq:adiabatic_cs_nonadiabatic_pressure} are obvious, since $p=w\rho$ with $w$ constant implies the fluid is barotropic $p=p(\rho)$.
Explicitly, differentiating the relation \eqref{eq:constant_eos_w} with respect to $y$ and $b$ respectively we get
\begin{align}
    y^2F_{yy}=t\frac{1}{w}\big(yF_y-byF_{by}\big), \qquad b^2 F_{bb}=w\big(bF_b-ybF_{yb}\big),
\end{align}
and plugging this into \eqref{eq:adiabatic_cs_nonadiabatic_pressure} we verify
{\small\begin{align}
\begin{aligned}
    &c_a^2=\frac{\partial p}{\partial \rho}_{|\sigma}=\frac{-y^2F_{yy}b^2F_{bb}+(yF_y-byF_{by})^2}{y^2F_{yy}(yF_y-bF_b)} \equiv w \\[3pt]
    &\Gamma:=\frac{\partial p}{\partial \sigma}_{\mid\rho}= \frac{b^2F_{bb}y^2F_{yy}-(yF_y-byF_{by})(bF_b-byF_{yb})}{\tfrac{1}{by}y^2F_{yy}(yF_y-bF_b)}\equiv 0.
\end{aligned}
\end{align}}



%========================================================================
\subsection{Fluids with time depdendent EoS parameter $w=w(a)$}
%=======================================================================
%
Despite reducing to the barotropic case, the solution for constant $w$ is nonetheless useful since a simple generalization is to consider a time-varying EoS parameter $w=w(a)$ with the scale factor $a$.
One might think the PDE is still formally solved by the above expression
\begin{align}
    F(y,b)=f(y\,b^{-w(a)})\, b^{(1+w(a))},
\end{align}
but this is NOT the case, since $y$ and $b$ ultimately depend on $a$ through the background equations.
We might still try the solution above, and then use the residual freedom in the form of the function $f$ to impose the other background equations, i.e. continuity and first Friedmann equation, giving consistency constraints on $f$ and its derivatives for the equations to hold.

It turns out it is more convenient to first solve for the evolution of the background variables $b(\tau),y(\tau)$ from the conservation of the Noether number-density and entropy currents \eqref{eq:conservation_noether_current_background_eft}.
The former \eqref{eq:conservation_noether_current_background_eft_number} implies that $b\propto a^3$ at the background level, so that $w=w(a)$ can be equivalently written as $w=w(b)$ and we are naturally lead to solving an equation of state with $w=w(b)$.
We refer to Section \ref{sec:bkg_evolution_fluid_EFT} for the treatment.


%=======================================================================
\subsection{Fluids with arbitrary EoS parameter $w(b,y)$}\label{sec:fluids_arbitrary_eos}
%==========================================================================

More generally, we can solve by characteristics the EoS for arbitrary parameters $w=w(b,\gamma)$
\begin{align}
    \frac{F-F_b}{yF_y-F}=w(b,y)\quad\Rightarrow\quad (1+w)F=w\, yF_y+bF_b.
\end{align}
This is just a linear transport equation, whose solution can be written for arbitrary $w$.
Choosing variables $\gamma:=\log (y/y_0)$ and $\beta:=\log (b/b_0)$, the characteristics are
{\small
\begin{align}\begin{aligned}
    \dot{\beta}=1,\qquad
    \dot{\gamma}=w(\gamma,\beta),\qquad
    \dot{z}=(1+w)z\,.
\end{aligned}\end{align}}
It is worth noticing that for $w=w(b)$ or $w=w(y)$ we could draw some further conclusions (e.g. on $c_a^2$ or $\Gamma$) from the sole equation of state, while for completely general $w(y,b)$ we cannot say anything more.
To proceed further we would need to couple the EoS with the background equations as in Section \ref{sec:bkg_evolution_fluid_EFT}.

We can use recursively the EoS to get relations between second dervatives of $F$ and simplify the expression for $c_a^2$.
Indeed, differentiating the equation and simplifying we get
{\small
\begin{align}\begin{aligned}
    &b^2F_{bb}= w \Big[(yF_y-byF_{by})-\left(1+w+\tfrac{bw_b}{w}\right)\tfrac{1}{w}(F-bF_b)\Big],\\
    &y^2F_{yy}=\tfrac{1}{w}\Big[(yF_y-byF_{by})-\tfrac{yw_y}{w}(F-bF_b)\Big],\\
\end{aligned}\end{align}}
We use these relations and the EoS $F-bF_b=w(yF_y-F)$ to simplify numerator and denominator of the expression for $c_a^2$
{\small
\begin{align}\begin{aligned}
    &-b^2F_{bb}y^2F_{yy}+(yF_y-byF_{by})^2\\
    &\quad=\tfrac{1}{w}(F-bF_b)\Big[(yF_y-byF_{by})\left(1+w+\tfrac{bw_b}{w}+\tfrac{yw_y}{w}\right)-\tfrac{yw_y}{w}\left(1+w+\tfrac{bw_b}{w}\right)(F-bF_b)\Big]\\
    &y^2F_{yy}(yF_y-bF_b)\\
    &\quad=y^2F_{yy}\tfrac{1+w}{w}(F-bF_b)=
    \tfrac{1}{w}\Big[(yF_y-byF_{by})-\tfrac{yw_y}{w}(F-bF_b)\Big] \tfrac{1+w}{w}(F-bF_b).
\end{aligned}\end{align}}
The expression for the speed of sound becomes 
{\small
\begin{align}\begin{aligned}
    c_a^2:&=\frac{-b^2F_{bb}y^2F_{yy}+(yF_y-byF_{by})^2}{y^2F_{yy}(yF_y-bF_b)}\\
    &=\frac{\cancel{\tfrac{1}{w}(F-bF_b)}\Big[(yF_y-byF_{by})\left(1+w+\tfrac{bw_b}{w}+\tfrac{yw_y}{w}\right)-\tfrac{yw_y}{w}\left(1+w+\tfrac{bw_b}{w}\right)(F-bF_b)\Big]}{\cancel{\tfrac{1}{w}(F-bF_b)}\tfrac{1+w}{w}\Big[(yF_y-byF_{by})-\tfrac{yw_y}{w}(F-bF_b)\Big] }
\end{aligned}\end{align}}
Rearranging and doing some further simplifications
{\small
\begin{align}\begin{aligned}
    (1+w)\left(\frac{c_a^2}{w}-1\right)
    &=\frac{\Big[(yF_y-byF_{by})\left(\tfrac{bw_b}{w}+\tfrac{yw_y}{w}\right)-\tfrac{yw_y}{w}\tfrac{bw_b}{w}(F-bF_b)\Big]}{\Big[(yF_y-byF_{by})-\tfrac{yw_y}{w}(F-bF_b)\Big]}\\[3pt]
    &=\left(\tfrac{bw_b}{w}+\tfrac{yw_y}{w}\right)+\left(\tfrac{yw_y}{w}\right)^2 \frac{(F-bF_b)}{(yF_y-byF_{by})-\tfrac{yw_y}{w}(F-bF_b)}
\end{aligned}\end{align}}
Finally, using logarithmic variables $\beta=\log b$ and $\gamma=\log y$ this becomes
{\small\begin{align}\begin{aligned}
\frac{d\log P}{d\log y}&= \frac{yF_y-byF_{by}}{F-bF_b}= \tfrac{w_\gamma}{w} + \frac{\left(\tfrac{w_\gamma}{w}\right)^2}{(1+w)\left(\frac{c_a^2}{w}-1\right)-\left(\tfrac{w_\beta}{w}+\tfrac{w_\gamma}{w}\right)}\\[3pt]
&= \frac{\left(\tfrac{w_\gamma}{w}\right)\Big[(1+w)\left(\frac{c_a^2}{w}-1\right)-\tfrac{w_\beta}{w}\Big]}{(1+w)\left(\frac{c_a^2}{w}-1\right)-\left(\tfrac{w_\beta}{w}+\tfrac{w_\gamma}{w}\right)}
\end{aligned}\end{align}}
Note the RHS is completely specified by $w$, we can thus integrate to get for a function $C=C(\beta)$ only
{\small\begin{align}\begin{aligned}
\log P&= C(\beta)+\int d\gamma\frac{\left(\tfrac{w_\gamma}{w}\right)\Big[(1+w)\left(\frac{c_a^2}{w}-1\right)-\tfrac{w_\beta}{w}\Big]}{(1+w)\left(\frac{c_a^2}{w}-1\right)-\left(\tfrac{w_\beta}{w}+\tfrac{w_\gamma}{w}\right)}
\end{aligned}\end{align}}
In turn we integrate $P=F-bF_b$ the ODE with known source given by $w$
{\small\begin{align}\begin{aligned}
F_\beta= F- e^{C(\beta)}\exp\left(\int d\gamma\frac{\left(\tfrac{w_\gamma}{w}\right)\Big[(1+w)\left(\frac{c_a^2}{w}-1\right)-\tfrac{w_\beta}{w}\Big]}{(1+w)\left(\frac{c_a^2}{w}-1\right)-\left(\tfrac{w_\beta}{w}+\tfrac{w_\gamma}{w}\right)}\right)
\end{aligned}\end{align}}

----------------------------------------------
In the case $c_a^2=1$ the expression becomes simpler 
{\small\begin{align}\begin{aligned}
\frac{d\log P}{d\log y}
&= \frac{\left(\tfrac{w_\gamma}{w}\right)\Big[\frac{1-w^2}{w}-\tfrac{w_\beta}{w}\Big]}{\frac{1-w^2}{w}-\left(\tfrac{w_\beta}{w}+\tfrac{w_\gamma}{w}\right)}= \left(\tfrac{w_\gamma}{w}\right)
\frac{\Big(1-w^2-w_\beta\Big)}{1-w^2-w_\beta-w_\gamma}
=  \left(\tfrac{w_\gamma}{w}\right) + 
\frac{\left(\tfrac{w_\gamma}{w}\right)^2}{\tfrac{1-w^2}{w}-\tfrac{w_\beta}{w}-\tfrac{w_\gamma}{w}}.
\end{aligned}\end{align}}

In the case of separable variables $w=u(\beta)v(\gamma)$ and $c_a^2=1$ if further simpliy to
{\small\begin{align}\begin{aligned}
\frac{d\log P}{d\log y}
&=  \left(\tfrac{v_\gamma}{v}\right) + 
\frac{\left(\tfrac{v_\gamma}{v}\right)^2}{\tfrac{1-u^2v^2}{uv}-\tfrac{u_\beta}{u}-\tfrac{v_\gamma}{v}}.
\end{aligned}\end{align}}

%-------------------------------------------------------------------------
\subsubsection{The cases $w=u(b)v(y)$}
%---------------------------------------------------------------------
%
For the moment we consider the case $w=u(b)v(y)$.
The characteristic for $\beta$ is trivial and is chosen as variable to parametrize the flow, while the solution for the $\gamma$ characteristic is
{\small
\begin{align}\begin{aligned}
    G(\gamma)-G(\gamma_0)=\int_{\beta_0}^\beta d\beta' u(\beta'), \quad \text{for}\quad G(\gamma):=\int^\gamma_* d\gamma' \frac{1}{v(\gamma')} + C\,.
\end{aligned}\end{align}}
The general solution is given for an arbitrary function of one variable $f=f(z)$ by
{\small
\begin{align}\begin{aligned}
    F(\beta,\gamma)=f\big(\tilde{\gamma}(\gamma,\beta,\beta_0)\big) \exp\left(\int_{\beta_0}^\beta d\beta' 1 + u(\beta')\,v\big(\tilde{\gamma}(\gamma,\beta,\beta')\big)\right), 
\end{aligned}\end{align}}
where the \emph{inverse} characteristics is defined by 
{\small
\begin{align}\begin{aligned}
    \tilde{\gamma}(\gamma,\beta,\beta'):=G^{-1}\left(G(\gamma)-\int^\beta_{\beta'}d\beta''\, u(\beta'')\right).
\end{aligned}\end{align}}
It is tedious but straightforward to check\footnote{I did painfully check.} this is indeed a solution using the relation
{\small
\begin{align}\begin{aligned}
    &-\frac{1}{u(\beta)}\,\frac{d}{d\beta}\tilde{\gamma}(\gamma,\beta,\beta')= \frac{1}{u(\beta')}\,\frac{d}{d\beta'}\tilde{\gamma}(\gamma,\beta,\beta')= v(\gamma) \frac{d}{d\gamma}\tilde{\gamma}(\gamma,\beta,\beta')=v(\tilde{\gamma}(\gamma,\beta,\beta'))\\[3pt]
    &\frac{d}{d\beta}\int_{\beta_0}^\beta \left[ 1 + u(\beta')\,v\big(\tilde{\gamma}(\gamma,\beta,\beta')\big)\right]d\beta'= 1+u(\beta)\,v(\tilde{\gamma}(\gamma,\beta,\beta_0))
\end{aligned}\end{align}}
The following relations are also useful for the thermodynamic computations
{\small
\begin{align}\begin{aligned}
    &F_\beta=F\left[1-u(\beta)v\big(\tilde{\gamma}(\gamma,\beta,\beta_0)\big)\left(\frac{f'\big(\tilde{\gamma}(\gamma,\beta,\beta_0)\big)}{f\big(\tilde{\gamma}(\gamma,\beta,\beta_0)\big)}-1\right)\right]\,,\\
    &F_\gamma = F\left[1+\frac{v\big(\tilde{\gamma}(\gamma,\beta,\beta_0)\big)}{v(\gamma)}\left(\frac{f'\big(\tilde{\gamma}(\gamma,\beta,\beta_0)\big)}{f\big(\tilde{\gamma}(\gamma,\beta,\beta_0)\big)}-1\right)\right]\,.
\end{aligned}\end{align}}
In particular
{\small
\begin{align}\begin{aligned}
    &p=F-F_\beta= F_\beta=F\left[u(\beta)v\big(\tilde{\gamma}_0\big)\left(\frac{f'\big(\tilde{\gamma}_0\big)}{f\big(\tilde{\gamma}_0\big)}-1\right)\right]\,,\\
    &\rho=F_\gamma-F = F\left[\frac{v\big(\tilde{\gamma}_0\big)}{v(\gamma)}\left(\frac{f'\big(\tilde{\gamma}_0\big)}{f\big(\tilde{\gamma}_0\big)}-1\right)\right]\,.
\end{aligned}\end{align}}

\textcolor{red}{Finally we compute $F_\beta\gamma$ and we plug the epxression in the simplified form of $c_a^2=1$ obtained above.}\todotag{finish}


%-------------------------------------------------------------------------
\subsubsection{The cases $w=w(b)$ and $w=w(y)$}
%---------------------------------------------------------------------
%
The most relevant case is $w=w(b)$ since the conservation of the Nother current \eqref{eq:conservation_noether_current_background_eft_number} gives $b\propto a^3$ at the backgrund level, so that  $w=w(b)$ ultimately corresponds to $w=(a)$ up to perturbations.
The equation of state is again
\begin{equation}\label{eq:eos_pde_w_depending_b}
    \frac{F-F_b}{yF_y-F}=w(b)\quad\Rightarrow\quad \big(1+w(b)\big)F=w(b)\, yF_y+bF_b.
\end{equation}
Choosing variables $\gamma:=\log (y/y_0)$ and $\beta:=\log (b/b_0)$, the characteristics are
{\small
\begin{align}\label{eq:characteristics_eos_pde_w_depending_b}
\begin{aligned}
    \beta^\prime&=1,\\
    \gamma^\prime&=w(\beta),\\
    z^\prime&=\big(1+w(\beta)\big)\,z
\end{aligned}\end{align}
}
The solution of the characteristics are simply 
{\small\begin{align}\label{eq:characteristics_eos_pde_solution_w_depending_b}
\begin{aligned}
    &\gamma(\beta,\gamma_0)=\int_{\beta_0}^\beta w(\beta)\,\,d\beta\,\,+\gamma_0\quad\Rightarrow\quad \gamma_0=\gamma-\int_{b_0}^b \!\!d\log(b/b_0)\,\, w(b)\,,\\
    &z(\beta):=F(\gamma(\beta,\gamma_0),\beta)= f(\gamma_0)\,\exp\left(\int_{\beta_0}^\beta\, d\beta\,\,(1+w(\beta))\right)\,.
\end{aligned}\end{align}}
The general solution for $w=w(b)$ is then
\begin{align}\label{eq:eos_pde_solution_w_depending_b}
    F(y,b)=f\Big(y\, e^{-\int_{b_0}^bd\log b\, w(b)}\Big)\, e^{\int_{b_0}^bd\log b\, \big(1+w(b)\big)},
\end{align}
which reduces to the constant EoS solution if $w=\text{const}$.
For the record, the general solutions for $w=w(y)$ is similarly 
\begin{align}
    F(y,b)=f\Big(b\, e^{-\int_{y_0}^yd\log y\, \frac{1}{w(y)}}\Big)\exp\left(\int_{y_0}^bd\log y\,\Big(1+\tfrac{1}{w(y)}\Big)\right).
\end{align}


